%=====================================================================
\chapter*{Nomenklatur\markboth{Nomenklatur}{Nomenklatur}}
\addcontentsline{toc}{chapter}{Nomenklatur}
%=====================================================================

\todo{Mathematische Sachen einbinden}

\setlength\LTleft{0pt}
% ###
% \subsection*{Lateinische Buchstaben}
% \begin{longtable}{@{\extracolsep{\LTleft}}p{20.0mm}p{25.0mm}l}
% $A$ & $\text{m}^2$ & Fl�che\\
% $\vec{a}$ & $\text{m/s}^2$ & Beschleunigung\\
% $a$ & $\text{m/s}$ & Schallgeschwindigkeit\\
% $c_p$ & $\text{J}/\left(\text{kg}\,\text{K}\right)$ & 
%     spezifische isobare W�rmekapazit�t\\
% $c_v$ & $\text{J}/\left(\text{kg}\,\text{K}\right)$ & 
%     spezifische isochore W�rmekapazit�t\\
% $E$ & $\text{J}$ & innere Energie\\
% $e$ & $\text{J/kg}$ & spezifische innere Energie\\
% $g$ & $\text{m/s}^2$ & Erdbeschleunigung\\
% $H$ & $\text{J}$ & Enthalpie\\
% $h$ & $\text{J/kg}$ & spezifische Enthalpie\\
% $k$ & $\text{J/K}$ & Boltzmann-Konstante $\left(k = 1,3807 \cdot 10^{-23}\right)$\\
% $m$ & $\text{kg}$ & Masse\\
% $N_\text{A}$ & $\text{mol}^{-1}$ & Avogadro-Konstante $\left(N_\text{A} = 6,0021318 \cdot 10^{23}\right)$\\
% $\vec{n}$ & $\text{-}$ & nach au�en gerichteter Normaleneinheitsvektor\\
% $n$ & $\text{-}$ & Polytropenexponent\\
% $P$ & $\text{J/s}$ & Leistung\\
% $p$ & $\text{N/m}^2$ & Druck\\
% $Q$ & $\text{J}$ & W�rme\\
% $\dot{Q}$ & $\text{J/s}$ & W�rmestrom\\
% $q$ & $\text{J/kg}$ & bezogene W�rme\\
% $\dot{q}$ & $\text{J}/\left(\text{kg}\,\text{s}\right)$ & 
%     bezogener W�rmestrom\\
% $R$ & $\text{J}/\left(\text{kg}\,\text{K}\right)$ & 
%     spezielle Gaskonstante\\
% $R_\text{m}$ & $\text{J}/\left(\text{mol}\,\text{K}\right)$ &
%     universelle Gaskonstante $\left(R_\text{m} = 8,31441 \right)$\\
% $T$ & $\text{K}$ & Temperatur\\
% $t$ & $\text{s}$ & Zeit\\
% $u, v, w$ & $\text{m/s}$ & Geschwindigkeitskomponenten im kartesischen\\
% & & Koordinatensystem\\
% $\vec{V}$ & $\text{m/s}$ & Geschwindigkeitsvektor\\
% $V$ & $\text{m}^3$ & Volumen\\
% $v$ & $\text{m}^3/\text{kg}$ & spezifisches Volumen\\
% $W$ & $\text{J}$ & Arbeit\\
% $x, y, z$ & $\text{m}$ & kartesische Koordinaten\\
% $z$ & $\text{m}$ & H�he
% \end{longtable}

% \subsection*{Griechische Buchstaben}
% \begin{longtable}{@{\extracolsep{\LTleft}}p{20.0mm}p{25.0mm}l}
% $\eta$ & $\text{Pa\,s}$ & dynamische Viskosit�t (Scherviskosit�t)\\
% $\kappa$ & $\text{-}$ & Isentropenexponent\\
% $\nu$ & $\text{m}^2/\text{s}$ & kinematische Viskosit�t\\
% $\nu$ & $\text{-}$ & Polytropenverh�ltnis\\
% $\varrho$ & $\text{kg/m}^3$ & Dichte\\
% $\tau_\text{w}$ & $\text{N/m}^2$ & Wandschubspannung\\
% $\vec{\omega}$ & $\text{s}^{-1}$ & Winkelgeschwindigkeit
% \end{longtable}

% \subsection*{Tiefgestellte Indizes}
% \begin{longtable}{@{\extracolsep{\LTleft}}p{20.0mm}l}
% $0$ & Ruhegr��e\\
% $\text{n}$ & normal\\
% $\text{t}$ & tangential
% \end{longtable}

% \subsection*{Hochgestellte Indizes}
% \begin{longtable}{@{\extracolsep{\LTleft}}p{20.0mm}l}
% $\ast$ & kritischer Zustand (\textsc{Laval}-Zustand)
% \end{longtable}

% \subsection*{Dimensionslose Kennzahlen}
% \begin{math}
% \begin{array}{@{\extracolsep{\LTleft}}lcll}
% \Ma & := & \norm{\vec{V}}/a & \mbox{Mach-Zahl}\\
% \Ma^\ast & := & \norm{\vec{V}}/a^\ast & \mbox{kritische Mach-Zahl}\\
% \Rey & := & \norm{\vec{V}}L/\nu & \mbox{Reynolds-Zahl}
% \end{array}
% \end{math}

% \subsection*{Operatoren, Funktionen und Symbole}
% \begin{longtable}{@{\extracolsep{\LTleft}}p{20.0mm}l}
% $\abs{\bullet}$ & Betrag eines Skalars\\
% $\norm{\bullet}$ & Norm (L�nge) eines Vektors\\
% d & totales bzw.\ vollst�ndiges Dif"|ferential\\
% $\deter(\bullet)$ & Determinante\\
% D & materielles (substantielles) Dif"|ferential\\
% $\grad(\bullet)$ & Gradient einer tensoriellen Gr��e\\
% $\delta$ & unvollst�ndiges Dif"|ferential\\
% $\partial$ & partielles Dif"|ferential\\
% $\partial V$ & st�ckweise glatter Rand eines Volumens\\
% $\nabla$ & Nablaoperator\\
% $\cdot$ & Skalarprodukt (inneres Produkt)\\
% $\times$ & Kreuzprodukt (�u�eres Produkt)\\
% \end{longtable}
% ###

\subsection*{Symbole}

\begin{longtable}{@{}p{3cm}p{11cm}@{}}
$x_{0}$            & Originaldaten (sauberes Volumen) \\
$x_{t}$            & Verrauschter Zustand zum Zeitschritt $t$ \\
$x_{T}$            & Reines Gau�sches Rauschen \\
$t,\ T$            & Aktueller Zeitschritt und Gesamtzahl der Zeitschritte \\
$\beta_{t}$        & Noise Schedule (Varianz pro Schritt) \\
$\alpha_{t},\ \bar{\alpha}_{t}$ & Signalanteil und dessen kumulatives Produkt \\
$\epsilon$         & Gau�sches Rauschen, $\epsilon\sim\mathcal{N}(0,I)$ \\
$\epsilon_{\theta}$ & Vom Netz vorhergesagtes Rauschen \\
$\theta$           & Lernbare Modellparameter \\
$\mathcal{N}(\mu,\sigma^{2})$ & Normalverteilung \\
$v$                & Velocity (Vorhersageziel der v-Prediction) \\
$w$                & Guidance-St�rke (Classifier-Free Guidance) \\
$\gamma$           & Min-SNR-Gewichtungsparameter \\
$z$                & Latente Variable bzw. Rauschvektor \\
\end{longtable}


\subsection*{Abk�rzungen}

% Abstand zwischen Abk�rzung und Beschreibung; ggf. anpassen
\renewcommand{\aclabelfont}[1]{\makebox[22mm][l]{#1}}

\begin{acronym}[]  % leere Liste nach acronym, sonst funktioniert die neu definierte textbox nicht und die Labels sind fett gegruckt.
    \itemsep=1.5pt  % Anpassung des vertikalen Abstands
    \acro{ADNI}{Alzheimer's Disease Neuroimaging Initiative}
    \acro{ANN}{Artificial Neural Networks}
    \acro{BAPPS}{Berkeley-Adobe Perceptual Patch Similarity}
    \acro{BraTS}{Brain Tumor Segmentation}
    \acro{cGAN}{conditional Generative Adversarial Network}        
    \acro{CNN}{Convolutional Neural Network}
    \acro{CT}{Computertomographie}
    \acro{DDIM}{Diffusion Denoising Implicit Model}
    \acro{DDPM}{Diffusion Denoising Probabilistic Model}
    \acro{DGM}{Deep Generative Model}
    \acro{DSC}{Dice Similarity Coefficient}
    \acro{ELBO}{Evidence Lower Bound}
    \acro{EMA}{Exponential Moving Average}
    \acro{FID}{Fr�chet Inception Distance}
    \acro{FNN}{Feedforward Neural Network}
    \acro{FR}{full-reference}
    \acro{GAN}{Generative Adversarial Network}                
    \acro{HA-GAN}{Hierarchical Amortized Generative Adversarial Network}
    \acro{IQMs}{Image Quality Metrics}
    \acro{IS}{Inception Score}
    \acro{IXI}{Information eXtraction from Images}
    \acro{KI}{K�nstliche Intelligenz}
    \acro{LDM}{Latent Diffusion Model}
    \acro{LPIPS}{Learned Perceptual Image Patch Similarity}
    \acro{LSGAN}{Least Squares Generative Adversarial Networks}
    \acro{MDP}{Markov Decision Process}
    \acro{ML}{Machine Learning}
    \acro{MLP}{Multi-Layer Perceptron}
    \acro{MRT}{Magnetresonanztomographie}
    \acro{MSE}{Mean Squared Error}
    \acro{NIfTI}{Neuroimaging Informatics Technology Initiative}
    \acro{NIH}{National Institutes of Health}
    \acro{NLP}{Natural Language Processing}
    \acro{NR}{no-reference}
    \acro{OASIS}{Open Access Series of Imaging Studies}
    \acro{ODE}{Ordinary Differential Equation}
    \acro{PSNR}{Peak Signal-to-Noise Ratio}
    \acro{ReLU}{Rectified Linear Unit}
    \acro{RNN}{Recurrent Neural Network}
    \acro{RR}{reduced-reference}
    \acro{SSIM}{Structural Similarity Index Measure}
    \acro{UniPC}{Unified Predictor-Corrector}
    \acro{VAE}{Variational Autoencoder Network}
    \acro{VQ-GAN}{Vector Quantized Generative Adversarial Network}
    \acro{WDM}{Wavelet Diffusion Model}
\end{acronym}